\section{Reproducibility}\label{sec:repro}

\paragraph{Source repository.}
The reference implementation source is held in a private repository,
\texttt{roastercode/radfi}, pending coordinated public release
with this methodology paper. Source access may be granted on a
case-by-case basis to qualified academic researchers under a light
non-disclosure agreement; contact
\texttt{aurelien@hackers.camp}. The present paper is published in
the public BEAMFS repository, \texttt{roastercode/beamfs},
sub-directory \texttt{papers/2026-04-radfi-v1/}.

\paragraph{Build environment.}
The reference build environment is a Yocto-based qemu-aarch64 lab
(Linux kernel 7.0.1, distro \texttt{poky-hardened} 1.0, kernel
module out-of-tree). The RadFI module loads via
\texttt{modprobe radfi}, exposes its debugfs interface under
\texttt{/sys/kernel/debug/radfi/}, and registers two kprobes whose
presence is verifiable through
\texttt{/sys/kernel/debug/kprobes/list}.

\paragraph{Empirical reproduction.}
Both empirical results of \cref{sec:empirical} are reproducible from
the public BEAMFS source code at the tagged commit:

\begin{itemize}
\item BEAMFS v1 falsification: \texttt{roastercode/beamfs} tag
  \texttt{v1.0.0-paper-pending} (commit \texttt{32f590e}), with
  RadFI v0.1.0 module load and the parameters of
  \cref{ssec:falsification}.
\item BEAMFS v2 INLINE confirmation: \texttt{roastercode/beamfs}
  tag \texttt{v0.2.0-palier3-validated} (commit \texttt{c47f130}),
  with the \texttt{dd}-based corruption procedure of
  \cref{ssec:confirmation} and the conformance fixture canary block
  defined in the BEAMFS reference
  implementation~\cite{beamfs-impl}, on-disk format v4
  specification (\texttt{Documentation/format-v4.md} \S 11.5).
\end{itemize}

The fixture SHA-256 values are normative and reproducible across
x86\_64 and aarch64 builds (verified by \texttt{cmp(1)} silent on
the canary block raw bytes). The normative fixtures are: raw
4{,}096-byte block
\texttt{\seqsplit{fb8a3b9e2704ce3fc11a0d0a4139d36c916cb5e7230acde8d2fa00ace739a91c}};
3{,}824-byte user content via VFS alias
\texttt{\seqsplit{ced446d9bf5e6682a48e8782ce5ad33f575f08921451afaa127f26dc37515bab}}.

\paragraph{Build (paper).}
The present paper is built from LaTeX sources in
\texttt{papers/2026-04-radfi-v1/} using \texttt{lualatex} +
\texttt{biber} + IEEEtran 2-column 10pt A4. Section files are
modular under \texttt{sections/}; bibliography is in
\texttt{refs.bib}. The build command is
\texttt{latexmk -lualatex paper.tex}.

\paragraph{License and distribution.}
Source code is under the GPL-2.0 license. The present technical
report is under CC-BY-4.0. Citation by tag and commit hash is
preferred over DOI for the implementation, as commit hashes are
strict identifiers of the audited state; DOI references for the
present report are added at Zenodo upload time.
